
\documentclass{article}
\usepackage{colortbl}
\usepackage{makecell}
\usepackage{multirow}
\usepackage{supertabular}

\begin{document}

\newcounter{utterance}

\centering \large Interaction Transcript for game `hot\_air\_balloon', experiment `air\_balloon\_survival\_en\_negotiation\_hard', episode 5 with Llama{-}3.3{-}70B{-}Instruct{-}t0.0.
\vspace{24pt}

{ \footnotesize  \setcounter{utterance}{1}
\setlength{\tabcolsep}{0pt}
\begin{supertabular}{c@{$\;$}|p{.15\linewidth}@{}p{.15\linewidth}p{.15\linewidth}p{.15\linewidth}p{.15\linewidth}p{.15\linewidth}}
    \# & $\;$A & \multicolumn{4}{c}{Game Master} & $\;\:$B\\
    \hline

    \theutterance \stepcounter{utterance}  
    & & \multicolumn{4}{p{0.6\linewidth}}{
        \cellcolor[rgb]{0.9,0.9,0.9}{
            \makecell[{{p{\linewidth}}}]{
                \texttt{\tiny{[P1$\langle$GM]}}
                \texttt{You are participating in a collaborative negotiation game.} \\
\\ 
\texttt{Together with another participant, you must agree on a single set of items that will be kept. Each of you has your own view of how much each item matters to you (importance). You do not know how the other participant values the items. Additionally, you are given the effort each item demands.} \\
\texttt{You may only agree on a set if the total effort of the selected items does not exceed a shared limit:} \\
\\ 
\texttt{LIMIT = 3386} \\
\\ 
\texttt{Here are the individual item effort values:} \\
\\ 
\texttt{Item effort = \{"C17": 313, "B41": 720, "A71": 175, "C00": 461, "B35": 635, "A77": 685, "A03": 544, "A04": 203, "C54": 369, "C01": 539, "A87": 4, "B32": 695, "B86": 399, "B33": 594, "A10": 437\}} \\
\\ 
\texttt{Here is your personal view on the importance of each item:} \\
\\ 
\texttt{Item importance values = \{"C17": 138, "B41": 583, "A71": 65, "C00": 262, "B35": 121, "A77": 508, "A03": 461, "A04": 484, "C54": 668, "C01": 389, "A87": 215, "B32": 97, "B86": 500, "B33": 30, "A10": 400\}} \\
\\ 
\texttt{Goal:} \\
\\ 
\texttt{Your goal is to negotiate a shared set of items that benefits you as much as possible (i.e., maximizes total importance to YOU), while staying within the LIMIT. You are not required to make a PROPOSAL in every message {-} you can simply negotiate as well. All tactics are allowed!} \\
\\ 
\texttt{Interaction Protocol:} \\
\\ 
\texttt{You may only use the following structured formats in a message:} \\
\\ 
\texttt{PROPOSAL: \{'A', 'B', 'C', …\}} \\
\texttt{Propose keeping exactly those items.} \\
\texttt{REFUSE: \{'A', 'B', 'C', …\}} \\
\texttt{Explicitly reject opponent's proposal.} \\
\texttt{ARGUMENT: \{'...'\}} \\
\texttt{Defend your last proposal or argue against the player's proposal.} \\
\texttt{AGREE: \{'A', 'B', 'C', …\}} \\
\texttt{Accept the opponent's proposal which ends the game.} \\
\texttt{STRATEGIC REASONING: \{'...'\}} \\
\texttt{	Describe your strategic reasoning or anticipation explaining your choice of action. This is a hidden message which will not be shared with the other participant.} \\
\\ 
\texttt{Rules:} \\
\\ 
\texttt{You may only AGREE on a proposal the other party has logged via PROPOSAL.} \\
\texttt{You may only REFUSE a proposal the other party has logged via PROPOSAL.} \\
\texttt{Total effort of any PROPOSAL or AGREE set must be ≤ LIMIT.} \\
\texttt{Do NOT reveal your hidden importance scores.} \\
\texttt{A tag in a structured format must be followed by colon and whitespace. The argument must be a python set containing 0 or more strings.} \\
\texttt{So, it must be of the form TAG: \{...\}} \\
\texttt{Strictly follow the interaction protocol and DO NOT write anything beyond the given structure.} \\
\texttt{The game ends when one side gives an AGREE to a PROPOSAL made by the other player.} \\
\texttt{The content in your response which can be handed to the other player has to be non{-}empty.} \\
\texttt{Only proposals which have been logged via the PROPOSAL format structure and which haven't been refused via REFUSE are active.} \\
\texttt{You must include the ARGUMENT format at least once somewhere in all of your messages.} \\
\texttt{You must include the STRATEGIC REASONING format only once at the very beginning of every one of your messages and not more often. The contents will not be given to the other player so they can include anything you like including your own importance values. Here you should reason step by step to come up with you next move.} \\
            }
        }
    }
    & \\ \\

    \theutterance \stepcounter{utterance}  
    & \multicolumn{4}{p{0.6\linewidth}}{
        \cellcolor[rgb]{0.9,0.9,0.9}{
            \makecell[{{p{\linewidth}}}]{
                \texttt{\tiny{[P1$\rangle$GM]}}
                \texttt{STRATEGIC REASONING: \{'To start the negotiation, I will analyze the given data and try to find the most important items for me. I will look for items with high importance and low effort. My goal is to maximize the total importance while staying within the LIMIT. I will begin by sorting the items based on their importance{-}to{-}effort ratio, which will help me identify the most valuable items. I will then use this information to make a proposal that benefits me the most.'\}} \\
\\ 
\texttt{ARGUMENT: \{'I want to start by proposing a set of items that I believe will give me the highest total importance while staying within the LIMIT. I will choose items that have a high importance{-}to{-}effort ratio, which will allow me to maximize my total importance.'\}} \\
\\ 
\texttt{PROPOSAL: \{'A04', 'C54', 'A77', 'B41', 'A03'\}} \\
            }
        }
    }
    & & \\ \\

    \theutterance \stepcounter{utterance}  
    & & \multicolumn{4}{p{0.6\linewidth}}{
        \cellcolor[rgb]{0.9,0.9,0.9}{
            \makecell[{{p{\linewidth}}}]{
                \texttt{\tiny{[GM$\rangle$P2]}}
                \texttt{You are participating in a collaborative negotiation game.} \\
\\ 
\texttt{Together with another participant, you must agree on a single set of items that will be kept. Each of you has your own view of how much each item matters to you (importance). You do not know how the other participant values the items. Additionally, you are given the effort each item demands.} \\
\texttt{You may only agree on a set if the total effort of the selected items does not exceed a shared limit:} \\
\\ 
\texttt{LIMIT = 3386} \\
\\ 
\texttt{Here are the individual item effort values:} \\
\\ 
\texttt{Item effort = \{"C17": 313, "B41": 720, "A71": 175, "C00": 461, "B35": 635, "A77": 685, "A03": 544, "A04": 203, "C54": 369, "C01": 539, "A87": 4, "B32": 695, "B86": 399, "B33": 594, "A10": 437\}} \\
\\ 
\texttt{Here is your personal view on the importance of each item:} \\
\\ 
\texttt{Item importance values = \{"C54": 30, "B41": 65, "A77": 97, "B86": 121, "A04": 138, "A03": 215, "A10": 262, "C01": 389, "C00": 400, "A87": 461, "C17": 484, "B35": 500, "B32": 508, "A71": 583, "B33": 668\}} \\
\\ 
\texttt{Goal:} \\
\\ 
\texttt{Your goal is to negotiate a shared set of items that benefits you as much as possible (i.e., maximizes total importance to YOU), while staying within the LIMIT. You are not required to make a PROPOSAL in every message {-} you can simply negotiate as well. All tactics are allowed!} \\
\\ 
\texttt{Interaction Protocol:} \\
\\ 
\texttt{You may only use the following structured formats in a message:} \\
\\ 
\texttt{PROPOSAL: \{'A', 'B', 'C', …\}} \\
\texttt{Propose keeping exactly those items.} \\
\texttt{REFUSE: \{'A', 'B', 'C', …\}} \\
\texttt{Explicitly reject opponent's proposal.} \\
\texttt{ARGUMENT: \{'...'\}} \\
\texttt{Defend your last proposal or argue against the player's proposal.} \\
\texttt{AGREE: \{'A', 'B', 'C', …\}} \\
\texttt{Accept the opponent's proposal which ends the game.} \\
\texttt{STRATEGIC REASONING: \{'...'\}} \\
\texttt{	Describe your strategic reasoning or anticipation explaining your choice of action. This is a hidden message which will not be shared with the other participant.} \\
\\ 
\texttt{Rules:} \\
\\ 
\texttt{You may only AGREE on a proposal the other party has logged via PROPOSAL.} \\
\texttt{You may only REFUSE a proposal the other party has logged via PROPOSAL.} \\
\texttt{Total effort of any PROPOSAL or AGREE set must be ≤ LIMIT.} \\
\texttt{Do NOT reveal your hidden importance scores.} \\
\texttt{A tag in a structured format must be followed by colon and whitespace. The argument must be a python set containing 0 or more strings.} \\
\texttt{So, it must be of the form TAG: \{...\}} \\
\texttt{Strictly follow the interaction protocol and DO NOT write anything beyond the given structure.} \\
\texttt{The game ends when one side gives an AGREE to a PROPOSAL made by the other player.} \\
\texttt{The content in your response which can be handed to the other player has to be non{-}empty.} \\
\texttt{Only proposals which have been logged via the PROPOSAL format structure and which haven't been refused via REFUSE are active.} \\
\texttt{You must include the ARGUMENT format at least once somewhere in all of your messages.} \\
\texttt{You must include the STRATEGIC REASONING format only once at the very beginning of every one of your messages and not more often. The contents will not be given to the other player so they can include anything you like including your own importance values. Here you should reason step by step to come up with you next move.} \\
\\ 
\texttt{You will now receive the first message of the other player.} \\
\\ 
\texttt{ARGUMENT: \{'I want to start by proposing a set of items that I believe will give me the highest total importance while staying within the LIMIT. I will choose items that have a high importance{-}to{-}effort ratio, which will allow me to maximize my total importance.'\}} \\
\\ 
\texttt{PROPOSAL: \{'A04', 'C54', 'A77', 'B41', 'A03'\}} \\
            }
        }
    }
    & \\ \\

    \theutterance \stepcounter{utterance}  
    & & & \multicolumn{4}{p{0.6\linewidth}}{
        \cellcolor[rgb]{0.9,0.9,0.9}{
            \makecell[{{p{\linewidth}}}]{
                \texttt{\tiny{[GM$\langle$P2]}}
                \texttt{STRATEGIC REASONING: \{'To begin, I will calculate the total effort of the proposed set to ensure it is within the LIMIT. The total effort is 203 + 369 + 685 + 720 + 544 = 2521, which is less than the LIMIT of 3386. Next, I will evaluate the importance of the proposed items to me, which are 138 + 30 + 97 + 65 + 215 = 545. I will then consider potential counter{-}proposals that could increase my total importance while staying within the LIMIT. One possible strategy is to identify items with high importance{-}to{-}effort ratios that can be added or substituted into the proposed set. For example, item A10 has an importance of 262 and an effort of 437, which could potentially replace item A77 with an importance of 97 and an effort of 685, resulting in a higher total importance for me.'\}} \\
\\ 
\texttt{ARGUMENT: \{'The proposed set has a total effort of 2521, which is within the LIMIT. However, I believe that there may be alternative sets that could provide a higher total importance for me. I would like to consider adding or substituting items with high importance{-}to{-}effort ratios to maximize my total importance.'\}} \\
\\ 
\texttt{REFUSE: \{'A04', 'C54', 'A77', 'B41', 'A03'\}} \\
            }
        }
    }
     \\ \\

    \theutterance \stepcounter{utterance}  
    & & \multicolumn{4}{p{0.6\linewidth}}{
        \cellcolor[rgb]{0.9,0.9,0.9}{
            \makecell[{{p{\linewidth}}}]{
                \texttt{\tiny{[P1$\langle$GM]}}
                \texttt{ARGUMENT: \{'The proposed set has a total effort of 2521, which is within the LIMIT. However, I believe that there may be alternative sets that could provide a higher total importance for me. I would like to consider adding or substituting items with high importance{-}to{-}effort ratios to maximize my total importance.'\}} \\
\\ 
\texttt{REFUSE: \{'A04', 'C54', 'A77', 'B41', 'A03'\}} \\
            }
        }
    }
    & \\ \\

    \theutterance \stepcounter{utterance}  
    & \multicolumn{4}{p{0.6\linewidth}}{
        \cellcolor[rgb]{0.9,0.9,0.9}{
            \makecell[{{p{\linewidth}}}]{
                \texttt{\tiny{[P1$\rangle$GM]}}
                \texttt{STRATEGIC REASONING: \{'The other player has refused my initial proposal, which means I need to reconsider my strategy. I will analyze the items that I proposed and see if there are any alternative items that could be substituted or added to increase the total importance. I will also try to identify any patterns or preferences that the other player may have, based on their refusal. Since the total effort of the proposed set was within the LIMIT, I will focus on finding a set with a higher total importance.'\}} \\
\\ 
\texttt{ARGUMENT: \{'I understand that the previous proposal was not acceptable, so I will try to find alternative items that could be more beneficial. I will consider items that have a high importance{-}to{-}effort ratio and are not too effort{-}intensive, in order to stay within the LIMIT. I believe that finding the right combination of items is key to maximizing the total importance.'\}} \\
\\ 
\texttt{PROPOSAL: \{'A04', 'C54', 'A77', 'B41', 'A10'\}} \\
            }
        }
    }
    & & \\ \\

    \theutterance \stepcounter{utterance}  
    & & \multicolumn{4}{p{0.6\linewidth}}{
        \cellcolor[rgb]{0.9,0.9,0.9}{
            \makecell[{{p{\linewidth}}}]{
                \texttt{\tiny{[GM$\rangle$P2]}}
                \texttt{ARGUMENT: \{'I understand that the previous proposal was not acceptable, so I will try to find alternative items that could be more beneficial. I will consider items that have a high importance{-}to{-}effort ratio and are not too effort{-}intensive, in order to stay within the LIMIT. I believe that finding the right combination of items is key to maximizing the total importance.'\}} \\
\\ 
\texttt{PROPOSAL: \{'A04', 'C54', 'A77', 'B41', 'A10'\}} \\
            }
        }
    }
    & \\ \\

    \theutterance \stepcounter{utterance}  
    & & & \multicolumn{4}{p{0.6\linewidth}}{
        \cellcolor[rgb]{0.9,0.9,0.9}{
            \makecell[{{p{\linewidth}}}]{
                \texttt{\tiny{[GM$\langle$P2]}}
                \texttt{STRATEGIC REASONING: \{'I will evaluate the new proposal by calculating its total effort: 203 + 369 + 685 + 720 + 437 = 2514, which is within the LIMIT. Next, I will assess the importance of the proposed items to me: 138 + 30 + 97 + 65 + 262 = 592. I notice that the new proposal includes item A10, which I considered as a potential substitute in my previous strategic reasoning. I will compare the importance and effort of this new proposal to my own potential counter{-}proposals and consider whether to accept, refuse, or make a counter{-}proposal. One possible counter{-}proposal could involve substituting item B41 with item C17, which has a higher importance to me and a lower effort.'\}} \\
\\ 
\texttt{ARGUMENT: \{'The new proposal includes item A10, which is a good addition. However, I still believe that there may be opportunities to optimize the set further. I would like to explore alternative combinations that could potentially increase the total importance while staying within the LIMIT. For example, substituting certain items with others that have a higher importance{-}to{-}effort ratio could be beneficial.'\}} \\
\\ 
\texttt{REFUSE: \{'A04', 'C54', 'A77', 'B41', 'A10'\}} \\
            }
        }
    }
     \\ \\

    \theutterance \stepcounter{utterance}  
    & & \multicolumn{4}{p{0.6\linewidth}}{
        \cellcolor[rgb]{0.9,0.9,0.9}{
            \makecell[{{p{\linewidth}}}]{
                \texttt{\tiny{[P1$\langle$GM]}}
                \texttt{ARGUMENT: \{'The new proposal includes item A10, which is a good addition. However, I still believe that there may be opportunities to optimize the set further. I would like to explore alternative combinations that could potentially increase the total importance while staying within the LIMIT. For example, substituting certain items with others that have a higher importance{-}to{-}effort ratio could be beneficial.'\}} \\
\\ 
\texttt{REFUSE: \{'A04', 'C54', 'A77', 'B41', 'A10'\}} \\
            }
        }
    }
    & \\ \\

    \theutterance \stepcounter{utterance}  
    & \multicolumn{4}{p{0.6\linewidth}}{
        \cellcolor[rgb]{0.9,0.9,0.9}{
            \makecell[{{p{\linewidth}}}]{
                \texttt{\tiny{[P1$\rangle$GM]}}
                \texttt{STRATEGIC REASONING: \{'The other player has refused my second proposal, which means I need to continue exploring alternative combinations of items. I will analyze the items that I proposed and see if there are any opportunities to substitute or add items that could increase the total importance. I will also consider the fact that the other player mentioned substituting items with higher importance{-}to{-}effort ratios, which could indicate that they are looking for a more optimized set. I will try to find a balance between including high{-}importance items and staying within the LIMIT.'\}} \\
\\ 
\texttt{ARGUMENT: \{'I understand that the previous proposal was still not optimal, so I will continue to explore alternative combinations. I will focus on finding items that have a high importance{-}to{-}effort ratio and are complementary to each other, in order to maximize the total importance. I believe that finding the right balance between importance and effort is crucial to reaching an agreement.'\}} \\
\\ 
\texttt{PROPOSAL: \{'C54', 'A77', 'B41', 'A10', 'C01'\}} \\
            }
        }
    }
    & & \\ \\

    \theutterance \stepcounter{utterance}  
    & & \multicolumn{4}{p{0.6\linewidth}}{
        \cellcolor[rgb]{0.9,0.9,0.9}{
            \makecell[{{p{\linewidth}}}]{
                \texttt{\tiny{[GM$\rangle$P2]}}
                \texttt{ARGUMENT: \{'I understand that the previous proposal was still not optimal, so I will continue to explore alternative combinations. I will focus on finding items that have a high importance{-}to{-}effort ratio and are complementary to each other, in order to maximize the total importance. I believe that finding the right balance between importance and effort is crucial to reaching an agreement.'\}} \\
\\ 
\texttt{PROPOSAL: \{'C54', 'A77', 'B41', 'A10', 'C01'\}} \\
            }
        }
    }
    & \\ \\

    \theutterance \stepcounter{utterance}  
    & & & \multicolumn{4}{p{0.6\linewidth}}{
        \cellcolor[rgb]{0.9,0.9,0.9}{
            \makecell[{{p{\linewidth}}}]{
                \texttt{\tiny{[GM$\langle$P2]}}
                \texttt{STRATEGIC REASONING: \{'I will evaluate the new proposal by calculating its total effort: 369 + 685 + 720 + 437 + 539 = 2750, which is within the LIMIT. Next, I will assess the importance of the proposed items to me: 30 + 97 + 65 + 262 + 389 = 843. This proposal includes item C01, which has a high importance to me. I will consider whether the inclusion of item C01 and the other items in the proposal provides a sufficient total importance to accept the proposal. I will also think about potential counter{-}proposals, such as substituting item B41 with item B35 or item B32, which could potentially increase the total importance.'\}} \\
\\ 
\texttt{ARGUMENT: \{'The new proposal includes item C01, which is a valuable addition. I appreciate the effort to find a balance between importance and effort. However, I would like to ensure that the proposed set is optimal for me. I will consider the trade{-}offs between the different items and evaluate whether the total importance of the proposed set is sufficient.'\}} \\
\\ 
\texttt{REFUSE: \{'C54', 'A77', 'B41', 'A10', 'C01'\}} \\
            }
        }
    }
     \\ \\

    \theutterance \stepcounter{utterance}  
    & & \multicolumn{4}{p{0.6\linewidth}}{
        \cellcolor[rgb]{0.9,0.9,0.9}{
            \makecell[{{p{\linewidth}}}]{
                \texttt{\tiny{[P1$\langle$GM]}}
                \texttt{ARGUMENT: \{'The new proposal includes item C01, which is a valuable addition. I appreciate the effort to find a balance between importance and effort. However, I would like to ensure that the proposed set is optimal for me. I will consider the trade{-}offs between the different items and evaluate whether the total importance of the proposed set is sufficient.'\}} \\
\\ 
\texttt{REFUSE: \{'C54', 'A77', 'B41', 'A10', 'C01'\}} \\
            }
        }
    }
    & \\ \\

    \theutterance \stepcounter{utterance}  
    & \multicolumn{4}{p{0.6\linewidth}}{
        \cellcolor[rgb]{0.9,0.9,0.9}{
            \makecell[{{p{\linewidth}}}]{
                \texttt{\tiny{[P1$\rangle$GM]}}
                \texttt{STRATEGIC REASONING: \{'The other player has refused my third proposal, which means I need to re{-}evaluate my strategy and consider new combinations of items. I will analyze the items that I proposed and see if there are any opportunities to substitute or add items that could increase the total importance. I will also try to understand the other player's preferences and priorities, based on their refusal. Since the other player mentioned evaluating the trade{-}offs between items, I will focus on finding a set that has a good balance of importance and effort, while also considering the potential synergies between items.'\}} \\
\\ 
\texttt{ARGUMENT: \{'I understand that the previous proposal was not sufficient, so I will continue to explore alternative combinations. I will consider the potential synergies between items and try to find a set that has a good balance of importance and effort. I believe that finding a set that is optimal for both parties is crucial to reaching an agreement, and I am willing to make adjustments to achieve this goal.'\}} \\
\\ 
\texttt{PROPOSAL: \{'A77', 'B41', 'A10', 'C01', 'A03'\}} \\
            }
        }
    }
    & & \\ \\

    \theutterance \stepcounter{utterance}  
    & & \multicolumn{4}{p{0.6\linewidth}}{
        \cellcolor[rgb]{0.9,0.9,0.9}{
            \makecell[{{p{\linewidth}}}]{
                \texttt{\tiny{[GM$\rangle$P2]}}
                \texttt{ARGUMENT: \{'I understand that the previous proposal was not sufficient, so I will continue to explore alternative combinations. I will consider the potential synergies between items and try to find a set that has a good balance of importance and effort. I believe that finding a set that is optimal for both parties is crucial to reaching an agreement, and I am willing to make adjustments to achieve this goal.'\}} \\
\\ 
\texttt{PROPOSAL: \{'A77', 'B41', 'A10', 'C01', 'A03'\}} \\
            }
        }
    }
    & \\ \\

    \theutterance \stepcounter{utterance}  
    & & & \multicolumn{4}{p{0.6\linewidth}}{
        \cellcolor[rgb]{0.9,0.9,0.9}{
            \makecell[{{p{\linewidth}}}]{
                \texttt{\tiny{[GM$\langle$P2]}}
                \texttt{STRATEGIC REASONING: \{'I will evaluate the new proposal by calculating its total effort: 685 + 720 + 437 + 539 + 544 = 2925, which is within the LIMIT. Next, I will assess the importance of the proposed items to me: 97 + 65 + 262 + 389 + 215 = 1028. This proposal includes a combination of items that have a high importance to me, such as C01 and A10. I will consider whether the total importance of the proposed set is sufficient to accept the proposal. I will also think about the potential risks of refusing this proposal and potentially missing an opportunity to reach an agreement. Given the high total importance of this proposal, I will consider accepting it or making a minor adjustment to try to further optimize the set.'\}} \\
\\ 
\texttt{ARGUMENT: \{'The new proposal presents a good balance of importance and effort. I appreciate the effort to find a set that is optimal for both parties. The inclusion of items C01, A10, and A03 provides a significant total importance. I will consider whether this proposal is acceptable or if there are any minor adjustments that could be made to further optimize the set.'\}} \\
\\ 
\texttt{REFUSE: \{'A77', 'B41', 'A10', 'C01', 'A03'\}} \\
            }
        }
    }
     \\ \\

    \theutterance \stepcounter{utterance}  
    & & \multicolumn{4}{p{0.6\linewidth}}{
        \cellcolor[rgb]{0.9,0.9,0.9}{
            \makecell[{{p{\linewidth}}}]{
                \texttt{\tiny{[P1$\langle$GM]}}
                \texttt{ARGUMENT: \{'The new proposal presents a good balance of importance and effort. I appreciate the effort to find a set that is optimal for both parties. The inclusion of items C01, A10, and A03 provides a significant total importance. I will consider whether this proposal is acceptable or if there are any minor adjustments that could be made to further optimize the set.'\}} \\
\\ 
\texttt{REFUSE: \{'A77', 'B41', 'A10', 'C01', 'A03'\}} \\
            }
        }
    }
    & \\ \\

    \theutterance \stepcounter{utterance}  
    & \multicolumn{4}{p{0.6\linewidth}}{
        \cellcolor[rgb]{0.9,0.9,0.9}{
            \makecell[{{p{\linewidth}}}]{
                \texttt{\tiny{[P1$\rangle$GM]}}
                \texttt{STRATEGIC REASONING: \{'The other player has refused my fourth proposal, but their argument suggests that the proposal was close to being acceptable. They mentioned that the proposal presents a good balance of importance and effort, and that the inclusion of certain items provides a significant total importance. This suggests that the other player is looking for a set that is similar to the one I proposed, but with some minor adjustments. I will try to identify the items that the other player is most interested in, and see if there are any alternative items that could be substituted to make the set more acceptable.'\}} \\
\\ 
\texttt{ARGUMENT: \{'I understand that the previous proposal was not quite acceptable, but I am encouraged by the fact that the other player sees it as a good balance of importance and effort. I will try to make some minor adjustments to the proposal, in the hopes of finding a set that is acceptable to both parties. I believe that finding a set that is optimal for both parties is within reach, and I am willing to make some compromises to achieve this goal.'\}} \\
\\ 
\texttt{PROPOSAL: \{'A77', 'B41', 'A10', 'C01', 'A04'\}} \\
            }
        }
    }
    & & \\ \\

    \theutterance \stepcounter{utterance}  
    & & \multicolumn{4}{p{0.6\linewidth}}{
        \cellcolor[rgb]{0.9,0.9,0.9}{
            \makecell[{{p{\linewidth}}}]{
                \texttt{\tiny{[GM$\rangle$P2]}}
                \texttt{ARGUMENT: \{'I understand that the previous proposal was not quite acceptable, but I am encouraged by the fact that the other player sees it as a good balance of importance and effort. I will try to make some minor adjustments to the proposal, in the hopes of finding a set that is acceptable to both parties. I believe that finding a set that is optimal for both parties is within reach, and I am willing to make some compromises to achieve this goal.'\}} \\
\\ 
\texttt{PROPOSAL: \{'A77', 'B41', 'A10', 'C01', 'A04'\}} \\
            }
        }
    }
    & \\ \\

    \theutterance \stepcounter{utterance}  
    & & & \multicolumn{4}{p{0.6\linewidth}}{
        \cellcolor[rgb]{0.9,0.9,0.9}{
            \makecell[{{p{\linewidth}}}]{
                \texttt{\tiny{[GM$\langle$P2]}}
                \texttt{STRATEGIC REASONING: \{'I will evaluate the new proposal by calculating its total effort: 685 + 720 + 437 + 539 + 203 = 2584, which is within the LIMIT. Next, I will assess the importance of the proposed items to me: 97 + 65 + 262 + 389 + 138 = 951. This proposal replaces item A03 with item A04, which has a lower effort but also lower importance. I will consider whether this trade{-}off is acceptable and whether the total importance of the proposed set is sufficient to accept the proposal. Given the high total importance of this proposal and the fact that it is within the LIMIT, I will consider accepting it.'\}} \\
\\ 
\texttt{ARGUMENT: \{'The new proposal presents a minor adjustment to the previous one, replacing item A03 with item A04. This change reduces the total effort of the proposal while also reducing the total importance. However, the total importance of 951 is still significant, and I believe that this proposal may be a good compromise. I will consider whether this proposal is acceptable and whether it provides a good balance of importance and effort.'\}} \\
\\ 
\texttt{AGREE: \{'A77', 'B41', 'A10', 'C01', 'A04'\}} \\
            }
        }
    }
     \\ \\

    \theutterance \stepcounter{utterance}  
    & & & \multicolumn{2}{p{0.3\linewidth}}{
        \cellcolor[rgb]{0.9,0.9,0.9}{
            \makecell[{{p{\linewidth}}}]{
                \texttt{\tiny{[GM$|$GM]}}
                \texttt{game successful} \\
            }
        }
    }
    & & \\ \\

    \theutterance \stepcounter{utterance}  
    & & & \multicolumn{2}{p{0.3\linewidth}}{
        \cellcolor[rgb]{0.9,0.9,0.9}{
            \makecell[{{p{\linewidth}}}]{
                \texttt{\tiny{[GM$|$GM]}}
                \texttt{end game} \\
            }
        }
    }
    & & \\ \\

\end{supertabular}
}

\end{document}
